% ===================== RESTAURAR ESTILO ANTES DEL FRONTMATTER =====================
\pagestyle{corporativo}

% ===================== FRONTMATTER CON MEMBRETE FORZADO =====================
\begin{frontmatter}

	\title{\projecttitle}

	\author[inst1]{J.Arboleda\corref{cor1}}
	\ead{proyectos2@dmlsas.com}

	\author[inst1]{F. Navia}
	\ead{fernando.navia@dmlsas.com}

	\author[inst2]{H. Rosero}
	\ead{herminsul.rosero@dmlsas.com}

	\cortext[cor1]{Autor correspondiente}
	\address[inst1]{DML Ingenieros consultores, Cali, Colombia, 760025}
	\address[inst2]{Especialidad Mecánica y Procesos, Cali, Colombia}

	% ===================== ABSTRACT CON MEMBRETE GARANTIZADO =====================
	\begin{abstract}
		\thispagestyle{corporativo}
		El presente informe técnico describe el diseño conceptual del sistema de ventilación con filtración para un laboratorio de 320 m$^3$ de volumen efectivo. El objetivo consiste en garantizar 12 renovaciones de aire por hora y diluir los contaminantes generados en el interior, con impulsión de aire filtrado y descarga libre a la atmósfera a través de rejillas.

		La estrategia de diseño seleccionada emplea un ventilador axial mural (placa mural) Ø560 mm en PRFV de transmisión directa (sin ductos), uniforme con el montaje típico de la planta, que introduce aire filtrado al recinto, complementado con tres rejillas de descarga libre que permiten la salida ordenada del flujo. Se presentan los cálculos de caudal (2 260 CFM), potencia del ventilador (0.320 kW teórico a 165 Pa en el sitio, motor provisional de 0.75 HP), velocidad real en boca (4.33 m/s a Ø560 mm) y dimensiones de las rejillas de salida (353 mm $\times$ 336 mm cada una, velocidad de descarga 3.0 m/s). Los resultados se sustentan en ASHRAE 170, ASHRAE 62.1, NFPA 99, OSHA 29 CFR 1910.1450, WHO Laboratory Biosafety Manual y NIH DRM.
	\end{abstract}

	\begin{keyword}
		ventilación de laboratorios \sep filtración MERV \sep renovaciones de aire \sep CFD \sep ASHRAE 170 \sep calidad del aire interior
	\end{keyword}

\end{frontmatter}

% ===================== RESTAURAR ESTILO DESPUÉS DEL FRONTMATTER =====================
\pagestyle{corporativo}
\thispagestyle{corporativo}
